\documentclass{article}

\usepackage[final]{neurips_2026}

\usepackage[utf8]{inputenc}
\usepackage[T1]{fontenc}
\usepackage[french]{babel}
\usepackage{hyperref}
\usepackage{url}
\usepackage{booktabs}
\usepackage{amsfonts}
\usepackage{nicefrac}
\usepackage{microtype}
\usepackage{xcolor}
\usepackage{graphicx}
\usepackage{subcaption}
\usepackage{amsmath}
\usepackage[title]{appendix}
\renewcommand{\appendixname}{Annexe}

\title{Topic Modelling des professions de foi aux législatives\\(1973--1993)}

\author{%
  Arthur MASCRET \\
  ENSAE \\
  \texttt{arthur.mascret@ensae.fr}
}

\begin{document}

\maketitle

\begin{abstract}
Ce travail applique trois méthodes de Topic modelling LDA, NMF et BERTopic à un corpus de 20\,478 professions de foi des élections législatives françaises de 1973 à 1993. L'objectif est d'extraire les thèmes dominants du discours électoral et de projeter les partis politiques dans l'espace thématique ainsi obtenu, afin d'en dresser une cartographie comparative. Nos résultats montrent que BERTopic, grâce à ses représentations contextuelles issues de transformers, produit des topics plus distincts et mieux séparés que LDA et NMF, dont les espaces thématiques restent concentrés autour d'un petit nombre de topics dominants. Les trois modèles s'accordent néanmoins sur de grands thèmes transversaux : le travail, le clivage travail-capital, la question sociale, l'écologie. La projection des partis dans l'espace thématique révèle une certaine homogénéité du discours électoral, tout en faisant apparaître des différenciations entre candidats politiques sur des enjeux spécifiques.
\end{abstract}

\section{Introduction}

En politique, il est souvent coutume de critiquer sans vergogne son adversaire ou un parti d'opposition dès qu'il prend position. Pourtant, il arrive que des formations politiques apparemment antagonistes tiennent un discours similaire sur certains enjeux. L'objectif de ce projet est d'utiliser le \textit{topic modelling} pour extraire les thèmes dominants des professions de foi électorales des législatives, puis de projeter les partis politiques dans cet espace thématique afin d'en dresser une cartographie. Contrairement aux programmes nationaux, les professions de foi reflètent le discours de chaque candidat dans son contexte local, révélant ainsi la diversité interne des formations politiques.

Nous exploitons pour cela les professions de foi des élections législatives conservées depuis la Troisième République, dont la numérisation est faite via la plateforme Arkindex (Sciences Po). Notre étude porte sur la période 1973--1993, riche sur le plan politique car elle englobe l'Union de la gauche et le Programme commun (1973--1978), la première alternance de la V\textsuperscript{e} République (1981), la première cohabitation (1986--1988), ainsi que la montée du Front National comme force parlementaire (1993). De plus, elle voit également l'émergence du mouvement écologiste.

Nous comparons trois méthodes de \textit{topic modelling} : Latent Dirichlet Allocation (LDA), la Non-negative Matrix Factorization (NMF) et BERTopic. Ces approches reposent sur des représentations textuelles différentes : sac-de-mots probabiliste pour LDA, décomposition TF-IDF pour NMF, et plongements lexicaux pour BERTopic.

\section{Rapide tour de littérature}


Le topic modelling désigne un ensemble de techniques d'apprentissage automatique non supervisées permettant de découvrir des thèmes latents dans de grands corpus. Ces approches reposent sur l'hypothèse que chaque document est composé d'un mélange de sujets, et que chaque sujet se caractérise par une distribution de probabilités sur un vocabulaire donné.

Parmi les approches matricielles, la factorisation en matrices non-négatives (NMF), introduite par Lee \& Seung~\cite{lee1999}, est une alternative déterministe pour le topic modelling. Le NMF décompose la matrice documents-termes en deux matrices de rang faible à coefficients positifs, représentant respectivement les distributions documents-topics et topics-mots. Cette contrainte de non-négativité favorise une interprétation plus intuitive des composantes, chaque topic étant représenté par une combinaison additive de mots.

L'algorithme Latent Dirichlet Allocation (LDA), proposé par Blei et al.~\cite{blei2003}, est l'un des modèles phares du Topic Modelling. Il modélise les documents comme des distributions de Dirichlet sur les topics, et les topics comme des distributions de Dirichlet sur les mots.

Les réseaux de neurones avec l'avènement des architectures Transformer ont révolutionné le traitement automatique du langage naturel grâce au mécanisme d'attention permettant de capturer les dépendances à longue distance dans les textes. Les modèles pré-entraînés comme BERT produisent des représentations contextuelles qui encodent une riche sémantique textuelle. Ces avancées ont bénéficié au topic modelling avec l'émergence de modèles comme BERTopic~\cite{grootendorst2022,bertopic_algo}, qui combine les embeddings de transformers avec des techniques de réduction dimensionnelle et de clustering pour découvrir des topics sémantiquement cohérents.



\section{Données et corpus}

\subsection{Description du corpus}

Le corpus s'étend sur cinq années 1973, 1978, 1981, 1988 et 1993. Après filtrage et appariement entre les documents numérisés et les métadonnées du fichier archelect\_search.csv, nous obtenons 20\,478 documents contenant en moyenne 555 mots (voir Annexe~\ref{ann:stats}). Les métadonnées associées à chaque document contiennent le parti du candidat, le sexe, la tranche d'âge, la profession, le statut de sortant, et le département.


\begin{figure}[h]
  \centering
  \begin{subfigure}[b]{0.41\linewidth}
    \centering
    \includegraphics[width=\linewidth]{figure/figure_1.png}
    \caption{Nombre de documents par année électorale.}
    \label{fig:docs_annee}
  \end{subfigure}
  \hfill
  \begin{subfigure}[b]{0.57\linewidth}
    \centering
    \includegraphics[width=\linewidth]{figure/figure_2.png}
    \caption{Répartition des 25 premiers partis/listes par nombre de documents.}
    \label{fig:partis}
  \end{subfigure}
  \caption{Répartition des professions de foi par élection et par parti}
\end{figure}



Comme le montre la Figure~\ref{fig:docs_annee}, le nombre de documents varie fortement selon les scrutins, passant d'environ 3\,800 à plus de 5\,800 entre 1973 et 1993. Le corpus se concentre autour des grands partis traditionnels, suivis d'une longue traîne de petits partis (Figure~\ref{fig:partis}).




\subsection{Prétraitement}
Le prétraitement constitue l'une des étapes les plus importantes du projet : sans un filtrage rigoureux, les topics extraits seront bruités par des mots non porteurs de sens, rendant leur interprétation impossible.

\begin{enumerate}
  \item \textbf{Lemmatisation} : réduction des mots à leur forme lemme (travaillons → travailler ) via spaCy.

  \item \textbf{Suppression des accents} : normalisation NFD suivie de filtrage des caractères combinants, permettant d'harmoniser les textes selon leur qualité d'OCR.

  \item \textbf{Filtrage des mots vides} : application d'une liste de 160 mots vides français augmentée de termes spécifiques au domaine (noms de partis, formules rhétoriques récurrentes, noms de politiciens, mots allemands pour les documents alsaciens). Cette suppression des noms de partis (« communiste », « républicain », etc.) est cruciale pour notre objectif, car nous cherchons à identifier les thèmes de discussion eux-mêmes, indépendamment de leur porteur politique.

  \item \textbf{Filtrage par fréquence} : les termes apparaissant dans moins de 10 documents sont écartés, tandis que les termes présents dans plus de 95\% des documents sont également exclus.
\end{enumerate}



\section{Modèles}

\subsection{LDA (Latent Dirichlet Allocation)}
Le modèle LDA est configuré avec $K=10$ topics, un vectoriseur de type \textit{bag-of-words} (CountVectorizer, 1\,000 features maximum, \texttt{min\_df}=5, \texttt{max\_df}=0.95) et un apprentissage en ligne (\textit{online} EM) sur 10 itérations.  Le nombre de thèmes $K=10$ est fixé après comparaison de la cohérence $C_V$ pour $K \in \{5, 10, 15\}$.

\subsection{NMF (Non-negative Matrix Factorization)}
La NMF décompose la matrice TF-IDF (TfidfVectorizer, 2\,000 features, \texttt{min\_df}=50, \texttt{max\_df}=0.95) en deux matrices non-négatives : $W \in \mathbb{R}^{n \times K}_+$ (documents $\times$ topics) et $H \in \mathbb{R}^{K \times m}_+$ (topics $\times$ mots). La matrice $W$ est normalisée ligne à ligne afin d'obtenir des distributions thématiques comparables. Le topic dominant de chaque document correspond à la colonne de $W$ ayant la valeur maximale. La décomposition utilise l'algorithme de descente de gradient par coordonnées avec $K=10$, fixé après comparaison de la cohérence $C_V$ pour $K \in \{5, 10, 15\}$.

\subsection{BERTopic}
BERTopic combine trois composantes principales :

\paragraph{Encodage sémantique.}
Chaque document est encodé avec le modèle \texttt{paraphrase-multilingual-MiniLM-L12-v2} (384 dimensions), pré-entraîné sur plus de 50 langues par apprentissage contrastif.

\paragraph{Réduction de dimensionnalité.}
UMAP réduit les représentations de 384 dimensions à 10 dimensions en préservant la structure locale (10 voisins proches).

\paragraph{Clustering et extraction de thèmes.}
HDBSCAN (\texttt{min\_cluster\_size}=10) regroupe les documents en clusters denses. Les thèmes sont représentés par leurs termes via c-TF-IDF, une variante qui pondère les termes par leur spécificité au cluster plutôt qu'au corpus global. Le modèle produit initialement un grand nombre de topics, réduits à 10 thèmes principaux par fusion hiérarchique.

\paragraph{Calibration des paramètres.}
Le réglage des paramètres de BERTopic a été effectué par tâtonnement. Une première exécution avec les paramètres par défaut et \texttt{nr\_topics="auto"} qui permet une réduction du nombre de topics. L'analyse du dendrogramme hiérarchique (Figure~\ref{fig:bert_hierarchy}) suggérait la présence de 9 topics distincts plus un cluster d'outliers. Nous avons donc exploré plusieurs combinaisons de paramètres en faisant varier \texttt{n\_neighbors}, \texttt{n\_components} et \texttt{min\_cluster\_size} pour ajuster respectivement : la structure locale/globale capturée, la richesse des représentations, et la granularité des clusters (peu de groupes volumineux versus nombreux petits groupes spécifiques). Le Tableau~\ref{tab:bert_params} synthétise les résultats de cette exploration.

\begin{figure}[h]
\centering
\includegraphics[width=0.8\textwidth]{figure/dendro.png}
\caption{Dendrogramme hiérarchique des topics BERTopic suggérant 9 topics}
\label{fig:bert_hierarchy}
\end{figure}

\begin{table}[h]
\centering
\caption{Exploration des hyperparamètres BERTopic}
\label{tab:bert_params}
\begin{tabular}{ccccccc}
\hline
\texttt{test} & \texttt{neighbors} & \texttt{components} & \texttt{min\_cluster} & Nb topics & U\_mass & C\_V \\
\hline
V0 & default (15) & default (5) & default (10) & 44 (non fixé) & -2.41 & 0.46  \\
V1 & 10 & 5 & 20 & 23 (non fixé) & -2.23 & 0.48  \\
V2 & 20 & 5 & 20 & 33 (non fixé) & -2.41 & 0.47  \\
V3 & 25 & 10 & 25 & 15 (non fixé) & -2.55 & 0.45  \\
V4 & default & default & default & 9 + 1 (fixé) & -1.86 & 0.52  \\
V5 & 10 & 10 & 10 & 9 + 1 (fixé) & -1.79 & 0.56  \\
V6 & 10 & 10 & 15 & 9 + 1 (non fixé) & -2.00 & 0.51  \\
V7 & 20 & 10 & 15 & 9 + 1 (fixé) & -1.88 & 0.51  \\
V8 & 10 & 15 & 10 & 9 + 1 (fixé) & -1.93 & 0.57  \\
V9 & 5 & 15 & 10 & 9 + 1 (non fixé) & -2.28 & 0.50  \\
\hline
\end{tabular}
\end{table}


\newpage
\section{Topics}

Le Tableau~\ref{tab:topics} présente les $K=10$ topics extraits par chaque modèle. Les labels sont attribués manuellement par interprétation des 15 termes les plus représentatifs de chaque topic, afin de faciliter la compréhension des topics (voir Annexe~\ref{ann:topwords}).

\begin{table}[h]
  \caption{Comparaison des topics par modèles}
  \label{tab:topics}
  \centering
  \small
  \begin{tabular}{clll}
    \toprule
    \# & LDA & NMF & BERTopic \\
    \midrule
    0 & Travail et entreprise            & Socialisme et autogestion ouvrière       & Travailleurs, femmes et Etat \\
    1 & Majorité et république           & Taxe      & Écologie, nature et entente  \\
    2 & Travailleurs et gouvernement     &  Féminisme              & Capital \textit{vs} travail          \\
    3 & Front populaire                  & Protection nature           & Socialisme et capitalisme            \\
    4 & Environnement                    & Immigration et impôts nationale       & Europe, illusion et économie         \\
    5 & Société et économie              & Écologie, entente et génération          & Emploi, retraite et logement         \\
    6 & Lutte ouvrière   & Institutions républicaines               & Écologie et génération               \\
    7 & Social et changement politique   & Anticapitalisme & Loi et criminalité                   \\
    8 & Impôt et immigration             & Capitalisme      & Corpus alsacien (all.)               \\
    9 & Écologie                         & Réforme sociale et pouvoir d'achat       & ---                                  \\
    \bottomrule
  \end{tabular}
\end{table}



L'attribution des labels repose sur l'analyse de l'importance des mots clés au sein de chaque topic. Cependant, cette approche présente certaines limites : la richesse sémantique est parfois restreinte, et certains topics peuvent sembler similaires ou peu informatifs, ce qui peut conduire à des erreurs de labellisation. C'est pourquoi les labels ont été croisés avec ceux générés par un LLM, qui perçoit mieux les relations sémantiques (notamment les cooccurrences lexicales) entre les mots. Cette démarche vise à obtenir les meilleurs résultats possibles tout en limitant le biais de subjectivité auquel on est confronté face à des thèmes ambigus.

On peut remarquer une certaine convergence des modèles vers de grands sujets, comme le clivage travail-capital ou l'opposition socialisme-capitalisme. Le travail constitue naturellement un des thèmes phares du corpus et l'environnement est également très présent. On retrouve aussi des questions d'ordre plus international, avec notamment l'Europe et l'immigration.

BERT semble produire des topics plus contrastés et précis que ses homologues, bien que le thème de l'écologie apparaisse deux fois dans ses résultats. Pour les deux autres modèles, on observe davantage de chevauchements entre topics, en particulier sur les thématiques liées à la gauche.




\begin{figure}[h]
  \centering
  \includegraphics[width=\linewidth]{figure/projection_topic.png}
  \caption{Projection des topics LDA (PCA), NMF (PCA) et BERTopic (UMAP)}
  \label{fig:proj_topics}
\end{figure}

La Figure~\ref{fig:proj_topics} semble confirmer cela, la projection des topics dans leur espace latent révèle une séparation plus nette pour BERTopic. En revanche, pour LDA, deux topics se situent à une distance extrême des autres (au-delà de 30 000 unités), ce qui écrase visuellement les autres topics sur ces deux dimensions. Pour NMF, la situation est plus satisfaisante que pour LDA : les topics sont globalement bien espacés et distribués sur une échelle cohérente, bien que certains restent relativement proches les uns des autres. Ces observations sont cohérentes avec le fait que le chevauchement entre topics est plus faible avec BERTopic qu'avec LDA et NMF.




\section{Projection des partis dans l'espace thématique}


Pour projeter chaque parti/liste dans l'espace thématique, nous calculons son \textit{profil thématique moyen} : pour chaque document, nous disposons d'un vecteur de distribution sur les $K$ topics (issu de \texttt{lda.transform}, de $W$ normalisé pour NMF, ou de \texttt{approximate\_distribution} pour BERTopic) ; le profil d'un parti est la moyenne de ces vecteurs sur l'ensemble de ses documents.

La projection 2D est ensuite réalisée différemment selon le modèle. Pour LDA et NMF, les composantes thématiques ($H$ pour NMF, \texttt{components\_} pour LDA) sont d'abord projetées par PCA en 2D ; les partis sont ensuite projetés dans ce même espace via le produit matriciel $\text{profil} \times H$, transformé par la même PCA.

Pour BERTopic, la logique est différente car le modèle fournit des embeddings sémantiques pour chaque topic. Le vecteur représentant un parti est construit comme une combinaison linéaire de ces embeddings, pondérée par son profil thématique : $\text{profil} \times E_{\text{topics}}$, où $E_{\text{topics}}$ est la matrice des embeddings de topics. L'ensemble des points (topics et partis) est ensuite projeté conjointement en 2D par UMAP, ce qui permet à la géométrie de l'espace de refléter les proximités sémantiques globales.
\subsection{LDA}



\begin{figure}[h]
  \centering
  \begin{minipage}[b]{0.48\linewidth}
    \centering
    \includegraphics[width=\linewidth]{figure/LDB.png}
    \caption{Projection partis LDA}
    \label{fig:proj_parties_lda}
  \end{minipage}
  \hfill
  \begin{minipage}[b]{0.48\linewidth}
    \centering
    \includegraphics[width=\linewidth]{figure/LDA_zoom.png}
    \caption{Projection partis LDA (zoom sur topics)}
    \label{fig:proj_parties_lda_zoom}
  \end{minipage}
\end{figure}


La projection conjointe topics/partis pour le LDA (Figures~\ref{fig:proj_parties_lda} et~\ref{fig:proj_parties_lda_zoom}) met en évidence une distribution fortement déséquilibrée dans l'espace thématique. Seuls 4 topics sur 10 concentrent la quasi-totalité des 1 189 partis/listes électorales, formant des groupes relativement distincts et nombreux. Les 6 topics restants, bien que relativement proches dans l'espace, correspondent à des zones quasi-vides en partis (Figure~\ref{fig:proj_parties_lda_zoom}). Ce peut suggérer qu'un LDA contraint à 4 topics aurait été pertinent.

Mais cela révèle aussi une limite : on pourrait interpréter qu'à l'échelle des professions de foi législatives, les 1 189 listes recensées partagent en réalité un vocabulaire très générique. Les candidats, quelle que soit leur étiquette politique, mobilisent les mêmes grands registres. Mais on sait que cela n'est pas le cas. De plus, c'est la répartition des partis dans l'espace des topics qui nous intéresse, pour identifier quelles listes traitent quels sujets, et dans quelle proportion. Avec seulement 4 topics discriminants, la granularité est trop faible pour mener une analyse de ce type. Une piste d'amélioration serait peut-être d'enrichir la liste des mots vides.
\subsection{NMF}

\begin{figure}[h]
  \centering
  \begin{minipage}[b]{0.48\linewidth}
    \centering
    \includegraphics[width=\linewidth]{figure/NMF.png}
    \caption{Projection partis NMF}
    \label{fig:proj_parties_nmf}
  \end{minipage}
  \hfill
  \begin{minipage}[b]{0.48\linewidth}
    \centering
    \includegraphics[width=\linewidth]{figure/NMF_zoom.png}
    \caption{Projection partis NMF (zoom)}
    \label{fig:proj_parties_nmf_zoom}
  \end{minipage}
\end{figure}

Les résultats produits par la méthode NMF s'avèrent globalement supérieurs à ceux du LDA, même s'ils demeurent imparfaits. On constate en effet que le corpus reste concentré autour d'un petit nombre de topics prédominants, sans parvenir à une répartition véritablement équilibrée. Malgré tout, cette méthode marque un progrès sensible par rapport au LDA, notamment en ce qui concerne la distinction entre groupes politiques. Toutefois, le faible nombre de topics réellement discriminants limite la portée de l'analyse. Il reste difficile d'examiner les positionnements de chaque candidat ou parti, car la répartition reste trop homogène pour permettre au modèle de faire émerger davantage de nuances dans les textes.



\subsection{BERTopic}

\begin{figure}[h]
  \centering
  \includegraphics[width=0.8\linewidth]{figure/bert.png}
  \caption{Projection partis BERTopic}
  \label{fig:proj_parties_bert}
\end{figure}

BERTopic a réussi à découper les thématiques de façon plus fine que les méthodes précédentes, où trois à quatre topics monopolisaient l'essentiel du corpus. Les groupes étant mieux séparés et répartis, il devient plus facile d'explorer en détail les sujets évoqués par chaque parti (voir Annexe~\ref{ann:bertopic}).

Pour aller plus loin, on pourrait analyser la disparité d'un parti en l'affichant seul sur le graphique, ce qui permettrait d'observer son niveau d'éclatement (discours homogène ou hétérogène). On pourrait également comparer uniquement deux partis que l'on suppose fortement opposés, comme le RN et LFI aujourd'hui, afin d'identifier d'éventuels points de convergence entre certains de leurs membres. De plus, l'analyse pourrait être étendue en étudiant les discours selon d'autres variables, comme l'âge ou le sexe, et non plus seulement en fonction du parti. Dans le cadre des législatives, où de nombreux candidats sont sans affiliation partisane, il serait intéressant d'examiner si certains d'entre eux se rapprochent fortement de partis existants dans les sujets évoqués.

Le topic « Travailleurs, femmes et gouvernement » se situe au cœur du graphique : les autres topics gravitent autour de lui. Cela est cohérent, puisqu'il s'agit du topic regroupant le plus grand nombre de partis. Il constitue ainsi un noyau commun aux autres thématiques, ce qui engendre une certaine hétérogénéité en son sein. En effet, bien que des partis appartiennent au même topic, ils peuvent se révéler très éloignés les uns des autres, voire parfois plus proches de partis appartenant à des topics distincts.

Ce phénomène illustre deux hypothèses. D'une part, malgré les bonnes performances de BERTopic, une marge d'amélioration subsiste dans la modélisation, notamment via un ajustement des hyperparamètres ou un prétraitement plus abouti. D'autre part, certains discours politiques sont par nature transversaux ou ambigus, naviguant entre plusieurs registres sans appartenir strictement à un seul.




\section{Conclusion}

Ce projet a exploré les professions de foi législatives françaises de 1973 à 1993 à l'aide du Topic modelling, en comparant trois approches : LDA, NMF et BERTopic.

Sur le plan des topics extraits, une convergence se dégage autour de grands clivages structurants sur le travail, capital, socialisme et capitalisme, écologie. BERTopic se distingue nettement par la qualité de sa segmentation thématique : ses topics sont plus cohérents, mieux séparés, et ses scores de cohérence $C_V$ sont supérieurs à ceux de LDA et NMF. La représentation contextuelle des textes par embeddings de transformers permet de capturer des proximités sémantiques que les approches bag-of-words ou TF-IDF ne parviennent pas à saisir.

La projection des partis dans l'espace thématique fait apparaître des résultats contrastés selon les modèles. Avec LDA, seulement quatre topics concentrent la quasi-totalité des partis, limitant fortement la granularité de l'analyse. NMF améliore sensiblement la distribution, mais sans atteindre une séparation satisfaisante. BERTopic offre la meilleure répartition, rendant possible une lecture plus fine des positionnements partisans et de leur éventuelle transversalité.

Ces résultats ouvrent plusieurs pistes d'amélioration. Sur le plan méthodologique, un enrichissement des stops words, une exploration des hyperparamètres, ou le recours à un modèle d'embeddings entraîné spécifiquement sur des textes politiques français pourraient affiner les résultats. Sur le plan analytique, il serait pertinent d'étendre l'analyse à d'autres variables sociodémographiques (sexe, âge, profession), d'examiner l'évolution temporelle des topics entre 1973 et 1993, ou encore de caractériser les candidats sans étiquette partisane par leur proximité thématique avec les grands partis.

Au-delà de ces limites, ce travail confirme que les professions de foi législatives constituent un corpus riche pour l'analyse du discours politique. Le topic modelling, et BERTopic en particulier, offre des outils prometteurs pour cartographier l'espace des idées à travers le temps et l'espace électoral.


\newpage
\section{Références}

\begin{thebibliography}{9}

\bibitem{lee1999}
Lee, D. D. \& Seung, H. S. (1999).
\newblock Learning the parts of objects by non-negative matrix factorization.
\newblock \textit{Nature}, 401(6755), 788--791.

\bibitem{blei2003}
Blei, D. M., Ng, A. Y. \& Jordan, M. I. (2003).
\newblock Latent Dirichlet Allocation.
\newblock \textit{Journal of Machine Learning Research}, 3, 993--1022.

\bibitem{grootendorst2022}
Grootendorst, M. (2022).
\newblock BERTopic: Neural topic modeling with a class-based TF-IDF procedure.
\newblock \textit{arXiv:2203.05794}.

\bibitem{bertopic_algo}
Grootendorst, M.
\newblock \textit{BERTopic -- The Algorithm}.
\newblock \url{https://maartengr.github.io/BERTopic/algorithm/algorithm.html}

\bibitem{bertopic_eval}
CSS -- Sciences Po / École Polytechnique.
\newblock \textit{BERTopic Tutorial: Evaluate your topic model}.
\newblock \url{https://css-polytechnique.github.io/css-ipp-materials/pages/bertopic-tutorial.html\#sec-evaluate-your-topic-model}

\bibitem{kermorvant2025}
Kermorvant, C. (2025--2026).
\newblock \textit{Cours de Traitement Automatique du Langage Naturel}.
\newblock ENSAE Paris.

\end{thebibliography}
\newpage
\begin{appendices}

\section{Statistiques descriptives}
\label{ann:stats}

\begin{figure}[h]
  \centering
  \includegraphics[width=0.6\linewidth]{figure/annexe.png}
  \caption{Statistiques descriptives Corpus}
  \label{fig:annexe_stats}
\end{figure}


\section{Termes caractéristiques par topic}
\label{ann:topwords}

\begin{figure}[h]
  \centering
  \begin{subfigure}[b]{\linewidth}
    \centering
    \includegraphics[width=0.8\linewidth]{figure/figure_6.png}
    \caption{LDA -- top 10 termes par topic.}
  \end{subfigure}
  \vspace{0.5em}
  \begin{subfigure}[b]{\linewidth}
    \centering
    \includegraphics[width=0.8\linewidth]{figure/figure_7.png}
    \caption{NMF -- top 15 termes par topic.}
  \end{subfigure}
  \caption{Termes les plus représentatifs pour LDA et NMF}
  \label{fig:topwords_lda_nmf}
\end{figure}

\newpage
\newpage
\section{Exploration interactive BERTopic}
\label{ann:bertopic}

La figure ci-dessous illustre l'interface d'exploration des partis/listes dans l'espace thématique BERTopic. Chaque point représente un parti, positionné selon sa distribution thématique moyenne, et coloré par topic dominant.

\begin{figure}[h]
  \centering
  \includegraphics[width=0.9\linewidth]{figure/bert_exemple.png}
  \caption{Exploration des partis/listes dans l'espace thématique BERTopic.}
  \label{fig:bertopic_exemple}
\end{figure}



\end{appendices}

\end{document}
